\section{Build sections around the reader's questions}
\label{sec:sections}

Write with a particular reader in mind. Each section should answer a question
the paper has raised, and each paragraph should carry one step of that answer.
This gives coauthors a practical test for duplication. Combine sections that
answer the same question. If useful details interrupt the argument, move them
to the appendix. A target page count cannot make either decision.

\subsection{Use the introduction to state the argument}

Open on the concrete task and the failure of current methods or theory. Narrow
quickly to the exact failure. The reader should understand the tension before
seeing the solution.

Then state the central insight in ordinary language. The method name and
technical machinery can follow. A reader who does not yet know the notation
should still be able to state the change and its importance.

The problem statement should open the paper's design questions. A useful form
names the starting object, the scarce resource, the resource that remains
available, and the outcome to explain. The rest of the introduction can then
separate questions about the representation, data, algorithm, and evaluation
instead of presenting one method as the entire answer.

Present contributions as supported claims:

\begin{quote}
\textbf{Sample efficiency.} We compare the method with [baselines] under
[controlled setting] and find [result] on [tasks or regimes].
\end{quote}

The heading names the claim, while the sentence gives the comparison and
evidence. ``We introduce a new algorithm,'' ``we run extensive experiments,''
and ``we provide theory'' name author activities. They report effort, not
findings.

For a paper that defines a constrained regime, the contribution sequence may
include the problem definition, a tested recipe, necessary conditions, a
characterization of the success and failure regimes of existing approaches,
and a new evaluation criterion. Each item still needs a result. List the
concrete findings before explaining their broader context.

Close the introduction with the paper's conceptual takeaway and its scope.
Do not repeat the same contribution in several formats. Choose one primary
signposting device and let the other elements add information.

\subsection{Place setup where readers need it}

Place each definition or assumption close to the first substantive result that
needs it. Keep standard preliminaries short. Spend space on a distinction that
later carries a claim, especially when two methods use similar notation but
optimize different objects.

A method section should explain the conceptual move before presenting every
implementation choice. Start from the limitation identified in the
introduction. Identify the changed computation and its purpose. Then give the
formal definition and practical variants.

A theory section should state the regime and result before a long derivation.
Readers need the theorem's purpose and target question. Keep the main proof
idea near the theorem when it carries scientific insight. Move routine algebra
and complete technical cases to the appendix. Never promote an intuition to
theorem language, and never omit a condition because it makes the statement
less elegant.

\subsection{State the question behind each experiment}

Begin the experimental section with the questions the experiments will answer.
Explain the distinction tested by each task and metric. Then justify the
baselines and state the controls. A comparison cannot support a clean mechanism
claim when the architecture, training objective, and initialization all change
at once.

The usual order is:

\begin{enumerate}[leftmargin=*]
    \item evaluation regimes and controlled comparison rules;
    \item headline performance;
    \item stress tests or deployment;
    \item mechanism studies;
    \item ablations of design decisions;
    \item efficiency, failures, and limitations.
\end{enumerate}

Change this order when the scientific question requires it. A mechanism paper
may need a controlled falsification before any broad benchmark. A paper about
sample cost should report environment interactions before wall-clock
convenience. The metric must match the claim.

State a result as a comparison. Name the metric and setting before offering an
explanation. Separate a performance comparison from a mechanism test.
``Method A outperforms Method B'' and ``component X causes the improvement''
are different claims and need different evidence.

Study each enabling condition separately. Vary the representation, task
coverage, update algorithm, and initial policy performance. Initial success
may not predict adaptability. Treat that difference as an experimental
question, not as an interpretation added after the results.

\subsection{Use related work to draw technical distinctions}

Group prior work by the technical choices that matter for the current paper.
For each group, explain its aim and the assumption that separates it from the
current paper. End with the point that remains unresolved. Cite the closest
alternatives directly. Broad claims such as ``few works consider this
setting'' require a careful search and a clear definition of the setting.

Treat a related method as evidence about a scientific distinction, not merely
as an opponent. Explain its complementarity with the current approach and name
the different resource or assumption behind each method. For readers outside
the subfield, describe the specialized method before relying on its acronym. A
citation supports that explanation; it does not replace it.

Related work should not become a second introduction. It should help the reader
place the paper's contribution accurately, including work that weakens the
novelty claim. A narrow, correct distinction is more credible than a broad
claim of being first. Cite distinctive language or framing that comes from
another paper. Do not disguise borrowed phrasing by rotating synonyms.

\subsection{Separate implications from limits in the discussion}

Use the discussion to interpret the evidence at the right level. State the
regimes in which the result should change practice or understanding. Then name
the limitation that most affects that conclusion. Untested environments may
narrow the result's breadth; a strong assumption may change where it applies.

The conclusion should not introduce a new contribution. It can state the
central result, its consequence, and the open question that now matters. A
generic recap weakens the ending. End on the most precise supported implication
or boundary.

\subsection{Make the appendix usable}

A long appendix benefits from a table of contents and an order designed for
readers. Put complete proofs and the details needed for reproduction there.
Practitioner guidance and qualitative failure analysis also belong in the
appendix when they help readers reuse the work.

The main paper must still support its central claim. Do not move a decisive
control or a condition that changes the theorem's meaning to the appendix
merely to save space. Move details that permit verification and reuse while
keeping the argument itself visible.
